\begin{abstract}
Online adaptation lets a speculative-decoding drafter keep up with the traffic it serves, but a better drafter does not always make serving faster. Whether speculation pays depends on load and on the target architecture, training the drafter on the serving GPU takes capacity away from requests, and some drafters cannot improve enough to ever pay off. We measure these effects on a dense (Qwen3-8B) and a mixture-of-experts (Qwen3-30B-A3B) target with EAGLE-3 drafters in vLLM, and build a controller that decides every second whether to speculate, how much of the GPU a co-located trainer may use, and when training should stop. The controller couples a calibrated profitability frontier, the break-even acceptance length as a function of concurrency, with per-domain learning curves: it trains only if the improved drafter can cross break-even and the gain expected over a forecast horizon repays the serving capacity spent on training. In a frozen evaluation of \nv{n.runs} paired runs (three scenarios, \nv{n.pol} policies, three repetitions), the best fixed policy differs between scenarios and every alternative falls at least \nv{alt.minshort}\% short of the per-scenario best in some scenario, while the controller stays within \nv{full.worstshort}\% of the best in all three. It raises saturated-phase throughput over matched static speculation by \nv{code.full.abs}\%, \nv{zh.full.abs}\%, and \nv{moe.full.abs}\%, turning speculation off when the drafter cannot reach break-even and adapting the drafter otherwise; on the MoE target the entire gain comes from online adaptation. Replayed in a frozen simulator on \nv{v1.n} unseen trace windows, the controller is never more than \nv{v1.short}\% below the best simple policy, and training pays back reliably only for bursts of about three minutes or longer.
\end{abstract}
